\section*{Chapter \ref{ch:g12:buffers-predominance} --- Buffers and Predominance Diagrams}

\begin{solution}{exo:g12:buffers-predominance:1}
\omterm{def:g9:acids-bases-ph:ph}{pH} 2: \ce{CH3COOH}. \omterm{def:g9:acids-bases-ph:ph}{pH} 4.76: both in equal amounts. \omterm{def:g9:acids-bases-ph:ph}{pH} 7: \ce{CH3COO-}.
\end{solution}

\begin{solution}{exo:g12:buffers-predominance:2}
\omterm{def:g9:acids-bases-ph:ph}{pH} 4: $1/(1 + 10^{0.76}) = 0.15$. \omterm{def:g9:acids-bases-ph:ph}{pH} 6: $1/(1 + 10^{-1.24}) = 0.95$.
\end{solution}

\begin{solution}{exo:g12:buffers-predominance:3}
Yellow; green (inside its range); blue.
\end{solution}

\begin{solution}{exo:g12:buffers-predominance:4}
At \omterm{def:g9:acids-bases-ph:ph}{pH} 7.4, below 9.26: \ce{NH4+}. At \omterm{def:g9:acids-bases-ph:ph}{pH} 11: \ce{NH3}.
\end{solution}

\begin{solution}{exo:g12:buffers-predominance:5}
An indicator is itself an \omterm{def:g12:ka-and-pka:bronsted}{acid--base couple}: in large amounts it would
change the \omterm{def:g9:acids-bases-ph:ph}{pH} it is meant to show.
\end{solution}

\begin{solution}{exo:g12:buffers-predominance:6}
$10^{\mathrm{pH} - 4.76}$: 0.1; 1; 10.
\end{solution}

\begin{solution}{exo:g12:buffers-predominance:7}
Phenolphthalein (8.0--10.0) or thymol blue (8.0--9.1): their ranges
contain 8.7. Methyl red changes between 4.2 and 6.3, far from 8.7: it
would have changed long before.
\end{solution}

\begin{solution}{exo:g12:buffers-predominance:8}
$\mathrm{pH} = 4.76 + \log(0.10/0.20) = 4.76 - 0.30 = 4.46$.
\end{solution}

\begin{solution}{exo:g12:buffers-predominance:9}
\omterm{def:g9:acids-bases-ph:ph}{pH} 1: the \omterm{def:g9:ions:ion}{cation} \ce{H3N+-CH2-COOH}, charge $+1$. \omterm{def:g9:acids-bases-ph:ph}{pH} 6: the zwitterion,
overall charge 0. \omterm{def:g9:acids-bases-ph:ph}{pH} 12: the \omterm{def:g9:ions:ion}{anion} \ce{H2N-CH2-COO-}, charge $-1$.
\end{solution}

\begin{solution}{exo:g12:buffers-predominance:10}
Water: from 7.0 to $-\log(\num{1.0e-4}/0.101) = 3.0$, a fall of 4 units.
Buffer: from 4.76 to $4.76 + \log(9.9/10.1) = 4.75$, a fall of 0.01.
\end{solution}

\begin{solution}{exo:g12:buffers-predominance:11}
\ce{NH4+}/\ce{NH3} ($pK_a = 9.26$). Ratio
$[\ce{NH3}]/[\ce{NH4+}] = 10^{9.5 - 9.26} = 1.7$.
\end{solution}

\begin{solution}{exo:g12:buffers-predominance:12}
The \omterm{def:g9:acids-bases-ph:ph}{pH} falls by one unit when the ratio reaches $1/10$:
$(10 - n)/(10 + n) = 0.1$, so $n = \qty{8.2}{mmol}$. A buffer holding more
of each form can absorb more acid or base for the same change of the
ratio: it has a larger capacity.
\end{solution}

\begin{solution}{exo:g12:buffers-predominance:13}
Ratio $10^{5.0 - 4.76} = 1.74$; the acid amounts to \qty{10}{mmol}, so
\qty{17.4}{mmol} of ethanoate, that is \qty{174}{mL} of the solution.
\end{solution}

\begin{solution}{exo:g12:buffers-predominance:14}
It has two \omterm{def:g12:ka-and-pka:bronsted}{acid--base couples}, with two $pK_a$, hence two changes of
colour: red at \omterm{def:g9:acids-bases-ph:ph}{pH} 1, yellow at \omterm{def:g9:acids-bases-ph:ph}{pH} 5, blue at \omterm{def:g9:acids-bases-ph:ph}{pH} 10.
\end{solution}

\begin{solution}{exo:g12:buffers-predominance:15}
\omterm{def:g10:concentration-and-dilution:dilution}{Dilution} divides both concentrations by 10: the ratio, and so the \omterm{def:g9:acids-bases-ph:ph}{pH}
(4.76), do not change. A \omterm{def:g12:ka-and-pka:strong-acid}{strong acid} diluted ten times rises by one \omterm{def:g9:acids-bases-ph:ph}{pH}
unit.
\end{solution}

\begin{solution}{pb:g12:buffers-predominance:1}
\textbf{1.} \ce{CO2 + 2H2O <=> HCO3- + H3O+}.

\textbf{2.} Dissolved carbon dioxide (with water) is the acid,
hydrogencarbonate the base.

\textbf{3.} The table value is for \qty{25}{\celsius} in pure water;
blood is at \qty{37}{\celsius} and full of other \omterm{def:g9:ions:ion}{ions}, which change the
constant (and physiologists count all the dissolved carbon dioxide).

\textbf{4.} At \qty{24}{mmol/L}, $0.024 \times 5.0 = \qty{0.12}{mol}$.

\textbf{5.} A \omterm{def:g9:acids-bases-ph:ph}{pH} axis cut at 6.1: \ce{CO2} below, \ce{HCO3-} above.

\textbf{6.} At 7.4, above 6.1: hydrogencarbonate.

\textbf{7.} Slightly basic: its \omterm{def:g9:acids-bases-ph:ph}{pH} is above 7.

\textbf{8.} $10^{7.4 - 6.1} = 10^{1.3} = 20$.

\textbf{9.} $24 / 20 = \qty{1.2}{mmol/L}$.

\textbf{10.} $10^{1.28} = 19$ and $10^{1.32} = 21$.

\textbf{11.} $20/21 = \qty{95}{\percent}$.

\textbf{12.} $10^{7.6 - 6.1} = 10^{1.5} = 32$.

\textbf{13.} The hydrogencarbonate \omterm{def:g9:ions:ion}{ion}:
\ce{HCO3- + H3O+ -> CO2 + 2H2O}.

\textbf{14.} $10^{7.1 - 6.1} = 10$.

\textbf{15.} From $24/20 = 1.2$ to $24/10 = \qty{2.4}{mmol/L}$: doubled.
The lungs breathe faster and deeper, to blow the extra carbon dioxide
out.

\textbf{16.} A \omterm{def:g12:ka-and-pka:strong-acid}{weak acid} and its base can each absorb what is added, and
the \omterm{def:g9:acids-bases-ph:ph}{pH} depends only on their ratio; a \omterm{def:g12:ka-and-pka:strong-acid}{strong acid} would only push the
\omterm{def:g9:acids-bases-ph:ph}{pH} one way.

\textbf{17.} Not at all: the ratio, and so the \omterm{def:g9:acids-bases-ph:ph}{pH}, stay the same.

\textbf{18.} About 20.
\end{solution}
